\section{Model setup}\label{sec:model_setup}

Having established the individual spatiotemporal signatures of FNS dynamics, we now define a standard computational framework to test these dynamics against the human behavioral data presented in \cref{chap:empirical_validations}. Because this thesis models human eye movements during visual search, we emulate the visual foraging task designed by our collaborators in \textcite{naher_primate_saccades}, which requires subjects to sequentially scan an array of dots to locate a hidden target.

\subsubsection*{Empirical protocol from Näher et al.}
In the human protocol established by \textcite{naher_primate_saccades}, trials commenced with a pre-search fixation period of 2500 to 3500~ms. Following a 100~ms gap, participants were presented with a search array of white dots, each subtending \qty{0.3}{\degree} of visual angle. Exactly 20\% of these dots functioned as true targets, and subjects were allotted a maximum search time of 5000~ms to successfully foveate a target. To investigate the influence of spatial structure, the study compared a REGULAR task---where dots were spaced identically to the centers of an invisible hexagonal tiling---with an IRREGULAR task, which randomly redistributed the same number of dots across the screen to break spatial regularity. Assuming a standard viewing distance of \qty{57.3}{\cm} (where $1^{\circ} \approx 1~\text{cm}$) and a monitor pixel density of \qty{36.22}{\px\per\cm}, spatial measurements can be directly translated between pixels and degrees using a scaling factor of \qty{36.22}{\px\per\degree}.

\subsubsection*{Computational model setup}
The empirical conditions from \textcite{naher_primate_saccades} are translated into numerical model parameters using units of spatial pixels and temporal seconds. The effective potential is constructed as a two-dimensional grid of Gaussian distributions, which act as quadratic potential wells. This simulated landscape consists of a $16 \times 16$ square grid of 256 potential wells, with 20\% of these wells randomly designated as true targets. Based on the average spatial distribution of eye movements during fixation, the standard deviation of the Gaussian wells is set to $\sigma_{\text{well}} = 3~\unit{\px}$. The average axial distance between wells is fixed at \qty{110}{\px}, which corresponds to a visual angle of approximately \qty{3.04}{\degree}. This spacing bridges the dense (\qty{2.3}{\degree}) and intermediate (\qty{4.5}{\degree}) spacings in REGULAR search conditions used in the human trials. To parameterise the spatial acuity scale of the dynamical-informative kernel $k(\cdot)$, $\sigma = \qty{80}{\px}$ was chosen to approximate human visual acuity. Visual acuity typically decays radially at a scale of roughly 2.5 degrees, which translates to 87.2 pixels assuming a viewing distance of 50 cm and a screen pixel density of 4 pixels per mm. Finally, the time constant $\tau = 1 \times 10^{-4}\,\unit{\second}$ was selected to align with the rhythmicity of human saccades. While the model yields an average inter-saccadic interval (ISI) of 360 ms, empirical data shows a saccade rate of 3.53 Hz, which corresponds to a 280 ms average ISI. Together, these elements define the 2D saliency map that acts as the explored effective potential. The full computational model whose solutions are studied throughout this work is:
\begin{align*}
    d\textbf{x}_t &= -\gamma c_{\alpha} \nabla V_{t}(\textbf{x})dt + \beta \textbf{p}_t dt + \gamma^{1/\alpha} dL_t^{\alpha},\\
    d\textbf{p}_t &= - \beta c_{\alpha} \nabla V_{t}(\textbf{x}) dt,\\
    \partial_t V_{t}(\textbf{x}) &= \frac{1}{\tau}k(\textbf{x} - \textbf{x}_t)
\end{align*}
Finally, the model parameters are calibrated to recover the distinct dynamical patterns analysed in \cref{chap:empirical_validations}. Thus, unless otherwise mentioned in this work, the baseline parameters considered for a standard experiment such as was conducted in \textcite{naher_primate_saccades}:
\begin{itemize}
    \item $\alpha = 1.2$ is selected to capture the strong search Lévy superdiffusion found in \cref{sec:search_superdiffusion} while maintaining numerical stability.
    \item $\sigma_{\text{k}} = 80~\unit{\px}$ in a kernel function $k(\textbf{x}) = \exp \{ \frac{-\|\textbf{x}\|^2}{2\sigma_{k}^2} \}$ scales the adaptive kernel to match the spatial limits of human visual acuity, as discussed above.
    \item $\tau = 2 \times 10^{-4}~\unit{\second^{-1}}$ is chosen to produce theta-rhythmic saccadic transitions, as derived in \cref{eq:theta_rhythm} and compared in \cref{sec:theta_rhythmicity_of_saccades}.
    \item $\beta = 1000~\unit{\px~ \second^{-1}}$ and $\gamma = 1000~\unit{\px^2~ \second^{-1}}$ ensures the recovery of gamma-rhythmic harmonic oscillations, placing the system well within the underdamped oscillatory regime for the \qty{3}{\px} potential wells. These are compared to data in \cref{sec:gamma_oscillations}
\end{itemize}

